\section{Discussion}
\label{sec:discussion}

\subsection{Which Configuration for Which Use Case}

The bakeoff results and the requirements matrix together answer the practitioner's
primary question: which configuration should I use?
We organise the recommendations by use case.

\paragraph{Default (no specific legal claim).}
Use \textbf{Geo+GeoSection} (config.~2).
It satisfies population equality, contiguity, and compactness (56\,\% improvement
over baseline), does not manipulate county lines or racial geography, and is
the best-documented configuration in the B-series.
The plan manifest is clean: only geographic signals are present.

\paragraph{Allen v.\ Milligan exposure (VRA Section 2, Gingles Prong 1).}
Use \textbf{Geo+VRASection} (config.~5), $w_\text{vra} = 0.40$.
This is the only configuration that achieves 2 minority-opportunity districts in
Alabama; the only configuration that achieves 2--3 in Georgia.
The 4.3\,\% compactness cost (EC premium) over GeoSection is legally defensible:
it falls within the range that courts have accepted for VRA compliance adjustments,
and the \miller{} racial-predominance test is satisfied by the $w_\text{vra} = 0.40$
parameterisation.

\paragraph{State-court partisan neutrality claim (LWV v.\ PA; Harper v.\ Hall).}
Use \textbf{Geo+AreaSection} (config.~3) to run the \emph{Rodden-null proof}:
demonstrate that the map is not a product of algorithmic partisan bias.
AreaSection reduces the proportionality gap in NC and GA by 7--8\,pp in our
estimates, not because the algorithm adds partisan signals but because it
prevents the urban-peeling geometry that creates artificially efficient Republican
configurations.
If the state constitution requires proportionality, ProportionalSection (B.12,
when implemented) is the direct tool---but it cannot be combined with VRASection
in the same run.

\paragraph{County preservation as a statutory criterion.}
Use \textbf{County+GeoSection} (config.~4) with $\alpha_\text{county}$ tuned to
the state's statutory standard.
For states like Wisconsin (strong county preservation tradition) or California
(explicit statutory criterion), a higher $\alpha$ reduces county splits at a
modest EC cost.
The Pareto frontier shows that a $50\,\%$ reduction in county splits typically
costs $5$--$10\,\%$ in normalised edge-cut.

\paragraph{Auditability and expert witness.}
Use \textbf{Geo+ApportionRegions} (config.~6) when the primary concern is
reproducibility and parameter-freedom.
The bisection tree is derived entirely from the prime factorisation of $k$;
there are no ``hidden'' parameter choices about which ratio to select.
Combined with the plan manifest, this provides the cleanest audit record.

\paragraph{Multi-chamber anti-stacking.}
Use \textbf{Geo+NestSection} (config.~8) in states where all three chambers
are drawn simultaneously and the NestSection compatibility score allows it.
The eleven states with spine score~0 (exact nesting possible) are the best
candidates; the remaining states require best-effort Mode~3 nesting.

\subsection{Limitations}

\subsubsection*{L1: Geographic Sorting Cannot Be Eliminated}

The most fundamental limitation is that geographic voter sorting produces
systematic Republican over-representation in competitive states regardless of
the algorithm.
All configurations show a mean proportionality gap of $-7$\,pp in the eight
competitive states studied across B.8--B.9.
AreaSection reduces this to approximately $-4$\,pp in the best case by preventing
the caterpillar pathology, but it cannot eliminate the underlying geographic
effect.
A state where Democrats are concentrated in three or four major cities and
Republicans are dispersed across the rural remainder will produce more Republican
seats than the vote-share would suggest under any geographically-neutral algorithm.
This is the central empirical finding of \citet{rodden2019}, and the B-series
programme confirms it across 44 states.

\subsubsection*{L2: VRASection Does Not Guarantee MM Districts}

VRASection's $2:5$ first-level split in Alabama concentrates minority population
in a geographic subregion; it does not guarantee that any specific district
within that subregion will be majority-minority.
The recursive GeoSection within the southern subregion must still achieve $\geq 50\,\%$
Black VAP in at least two districts.
B.14 confirms that this does occur (55.2\,\% and 52.8\,\% Black VAP in the two
southern districts), but the guarantee is empirical, not structural.
For states where the minority population is more dispersed, or where the $2:5$
subregion has Black VAP just above 50\,\%, the guarantee may not hold.
Users should run \texttt{redist analyze --types demographic} to verify MM
district counts after every VRASection run.

\subsubsection*{L3: Census Stability is Not Universal}

StabilitySection (B.15) shows that 33\,\% of states are \emph{unstable}---their
GeoSection natural ratio changes between census years.
The census-stability litigation argument cannot be made for these states.
Iowa is the canonical example: its natural ratio shifted from $1:3$ (2000) to
$2:2$ (2010) as suburban population dispersal changed the population-geographic
structure.
For Iowa-class states, the practitioner should present the map as reflecting
the \emph{current census year's geographic structure}, not a census-stable
property.

\subsubsection*{L4: ProportionalSection is Planned, Not Implemented}

The ProportionalSection (B.12) and ApportionRegions (B.11) components are
described in this synthesis but not yet implemented in the \texttt{redist} CLI.
Results for configs.~6 and 8 in the bakeoff tables are pending.
The mutual-exclusion architecture described in Section~\ref{sec:callais} is
fully implemented and will enforce the ProportionalSection mutex as soon as
B.12 is complete.

\subsubsection*{L5: County-Sticky Weights Require TIGER Download}

Full subdivision-stickiness (MCD, place, VTD levels) requires downloading
additional TIGER spatial files ($\sim$2\,GB for all 50 states) and performing
spatial joins.
County-level stickiness is available immediately (county FIPS is encoded in the
census-tract GEOID); finer levels require the B.10 Phase~2--3 implementation.

\subsection{Comparison with Existing Approaches}

\subsubsection*{Markov Chain Methods (GerryChain)}

GerryChain \citep{metric2018} and related Markov-chain Monte Carlo (MCMC) approaches
generate large ensembles---typically thousands of independently drawn plans---and test
whether an enacted map falls within the ensemble's distribution.
These methods are well-suited to null-hypothesis testing (``is the enacted map
unusual relative to the space of plans?'').
The toolbox is complementary: it generates a single deterministic plan optimising
a specified criterion, while GerryChain generates an MCMC ensemble of thousands of
plans to characterise the solution space.
A practitioner may use the two approaches together: run the toolbox to obtain the
certified-optimal plan, then use GerryChain to verify that the plan lies in the
interior of the ensemble (i.e., is not an outlier in any partisan or demographic dimension).
B.15's Census Stability Score provides a deterministic stability metric that
MCMC cannot easily compute (it would require running an ensemble for each of
three census years).

\subsubsection*{Integer Programming}

Integer programming (IP) approaches (e.g., \citet{gurnee2021}) can optimise
multiple objectives simultaneously but do not scale to large states.
For California (52 districts, $\sim$8,000 census tracts), IP is intractable.
The METIS-based toolbox operates in O($n \log k$) time and produces results
in 30--90 seconds per state on a single CPU.
The trade-off is that METIS produces locally optimal (not globally optimal)
solutions; the B.7 seed study shows this produces a narrow solution space
with $< 0.3$ seats of variation per state.

\subsubsection*{Short Burst Methods}

Short burst methods \citep{cannon2023} run many short Markov chains to
explore the redistricting solution space.
They are well-suited to understanding solution space geometry but share the
MCMC limitation that they are not deterministic: two independent runs of the
same method on the same state may produce different plans.
The toolbox's auditable determinism (plan manifest + hash chain) is a
structural advantage for litigation, where reproducibility is a requirement,
not a preference.

\subsection{The Algorithm as Evidence, Not Advocacy}

A final point deserves emphasis.
The toolbox is designed as an \emph{evidentiary} tool, not an advocacy tool.
It does not argue for a particular redistricting outcome.
It provides a set of analytically defensible configurations, each grounded in
a specific legal requirement, each producing a reproducible and auditable result.

A redistricting commission, a court-appointed special master, or an opposing
expert witness can all use the same toolbox and reach different conclusions---
because the conclusions depend on which legal requirements are weighted most
heavily, not on which algorithm is used.
VRASection and GeoSection will sometimes produce different seat counts (North
Carolina, Wisconsin) and sometimes not (Alabama with $0\mathrm{D}$ in both).
The algorithm does not choose; the legal argument and the factual record choose.
The toolbox provides the counterfactual analysis that makes the argument
\emph{auditable}.
